\subsection*{Independent second reading of the automorphism-orbit criterion}

All 317 lines of \texttt{finite\_checks/b1\_automorphism\_orbits.tex}
and all 278 lines of its finite certificate program were
independently read on 18 November 2026.  The retained
key outputs were read as well.  The outcome for B.145
is \textbf{PASS}, without a requested mathematical
correction.  This is a separate same-model-family check
by GPT-6 Astra Pro in Pro mode; a separate effort setting
is not exposed.  Code reading is distinguished from
executing that complete certificate.

For an unramified nontrivial character, the effective
divisor expansion and Riemann--Roch give exactly
$A_d(i)=q^{d+1-g}\zeta^{ic}A_{2g-2-d}(-i)$ in the
indicated range.  The common class sum vanishes,
including its constant term.  Thus reduction modulo
$p$ leaves degree at most $g-1$ and constant term one.
There is no ordinarity assumption.

The automorphism identity has the stated sign:
permuting divisors gives
$A_d(i)=\zeta^{ida}A_d(iu)$.  Raising a vanishing
factor to $p^j$ therefore sends its twist to
$ut-a$.  The affine fixed point is $a/(u-1)$;
every other orbit has length $m$.  The degree bound
excludes such an orbit of roots, simultaneously for
all nontrivial characters.  The argument allows
wild automorphisms, since it uses a permutation of
divisors and does not average over the automorphism
group.  The norm character after constant extension
preserves its eigenvalue and multiplies the marked
shift by the extension degree, giving the precise
prime-to-$\ell$ persistence assertion.

The example is a smooth plane quartic: its three
partial derivatives have no common projective zero.
The infinity line has divisor $4O$, giving the
canonical class.  The degree-two elliptic map and
the translation automorphism intertwine as stated.
At $R=(t,1)$, the elliptic doubling formula gives
$2R=(t+1,1)$ and $4R=(t+2,1)$, hence exact order
seven.  The elliptic group over $\mathbb F_{27}$
has order $28$ and splits as its rational
two-torsion and the cyclic group generated by $R$.
This proves the eigencharacter and marked shift
before the larger finite computation is used.

The certificate's two finite fields have the stated
irreducible defining polynomials.  Its degree-two
point loop picks precisely one element of each
nonrational conjugate pair and uses the elliptic
norm to compute the divisor character.  Together
with all unordered pairs of rational points, with
repetition, this is the complete degree-two
effective-divisor count.  The retained occupancy
totals are $28$ and $730=406+324$ as required.

Newton's identities give
$(X^2-3)^2(X^2+3)$ from the first three curve
counts; cubing the roots gives the displayed
$\mathbb F_{27}$ polynomial and class number.
The high character coefficients use the correct
canonical exponent $3$.  Reducing the degree-two
occupancy gives $-82\zeta^5$, and evaluation at
one gives the written polynomial $A(Z)$ up to a
multiple of $\Phi_7$.  The retained resultant and
multiplication determinant agree.  Independently,
its displayed reduction modulo three is nonzero
and not a multiple of the irreducible degree-six
$\Phi_7$, which proves the required unit assertion.
The minimal-prime claim is only at the stated
field: five is excluded by the actual base
Jacobian order, and three is the characteristic.

The character table gives unpointed elliptic values
$0$ at the four original rational points and
$1,2,4$ each eight times at the other points.
The automorphism cycles the latter three classes,
so all their marked shifts are nonzero and their
relative class numbers agree.  This verifies the
claimed twenty-four-marking family.

\medskip\noindent\textbf{An additional check sent to the author.}
The one twist not certified by the general
criterion is nevertheless good in this example.
At $t=1$, direct evaluation gives
\[
 A_1(Z)=760+8Z+8Z^2+216Z^3+8Z^4+216Z^5+216Z^6.
\]
Its reduction is $1+2Z+2Z^2+2Z^4$, nonzero of
degree below six.  An independent exact resultant
calculation gave
$121740744925904896$ for the relative class number
and $2304308819957527871488$ for the cover Jacobian
order.  This concerns the four original markings
after the stated extension to $\mathbb F_{27}$;
it is separate from the general nonzero-shift
criterion and was communicated for a separately
checked addition.

The quadratic-period identity also shows that
the primitive unit roots of this twist have order
four modulo three; the other twists multiply
them by nontrivial seventh roots and have order
twenty-eight.  These are tests for these explicit
covers under constant extension, and by themselves
are not a complete least-witness census.

\medskip\noindent\textbf{Final addon: separate targeted reading.}
The final 412-line proof's added statement clauses
and all added proof paragraphs were read on
18 November 2026, together with the expanded
twist and period certificate code and its retained
outputs.  The outcome is \textbf{PASS}.
The seven twists give exactly the three displayed
integer norms.  Resultant and multiplication-matrix
formulas use the same exact cyclotomic coefficients.
Changing an original marking of character zero
from the new marking of character one gives the
positive twist one in every degree.

The identity $U^2+U+2=0$ gives $\eta^2=-1$,
and the coefficient calculation is exactly
$(1-\eta\bar\zeta^{t-1}T)^2$ modulo three.
Its nonzero roots have order four at twist one
and order twenty-eight otherwise, because four
and seven are coprime.  Every primitive conjugate
has the same order.  Under every constant
extension the roots are raised to its degree;
roots with zero reduction cannot make evaluation
at one divisible by three.  This proves the
all-degree tests without extrapolation.

The exact base Jacobian orders $16$, $256$
and $18928$ exclude every odd-prime constant
pointed cyclic cover at the first two fields.
The third admits the checked degree-seven
twist-one witness, with no smaller admissible
prime.  Thus the field-first minimum $(3,7)$
has precisely the stated scope.  This event
records complete reading of the expanded
certificate, not independent execution of
that complete program.
